% ===== §9 Formal Expression: The Translation Layer =====
\section{Formal Expression: The Translation Layer}
\label{p20:sec:formal}

\noindent
The preceding sections stated, at the points where they needed them, the formal results that underwrite the ethics, and pointed here for their full form. This section keeps that promise. It is a translation layer and nothing more: it does not build a formalism, it does not prove the results afresh, and it claims no mathematical novelty. Every result below is established in the formal companion to this series \parencite{wan2026braid}; what is done here is to set each beside the water-virtue it expresses, so that the reader may see that the ethical claims of the paper and the theorems of the companion are the same claims in two registers. The serving position is strict: where the ethics and the formalism diverge, it is the ethics that governs, and the formalism is retained only insofar as it expresses something the ethics independently holds. The truth-semantics that licenses reading these theorems as claims about the good, rather than as claims about a model, is the thin relational realism declared in the survey (\xref{p20:ssub:review-metaethics}): the water-good is a structural property of the flow, and a description of that structure is a description of the good.

\subsection{Non-contention as non-possessive surplus}
\label{p20:sub:formal-noncon}

The companion models a relation's value as carried by a flow, and its generated surplus as the holonomy of that flow around a closed relational cycle --- the value a complete circuit returns over and above what entered it, $\holo(\gamma) = \oint_\gamma u \cdot d\ell$. Its founding theorem is that the appropriative part of the driving force, being the gradient of a scalar potential, contributes nothing to this holonomy: a purely appropriative flow $u = -\nabla\phi$ has $\holo(\gamma) = 0$ around every closed loop \parencite{wan2026braid}. This is the exact form of the claim of \xref{p20:sub:noncon-bushi}. Whatever in a relation reduces to a single scalar ordering --- possession, ranking, equivalent exchange, all of which are gradients along a ``more or less'' --- generates no surplus; only the non-gradient, circulating part generates anything. Non-contention is, formally, the relation's reliance on the circulating rather than the gradient part of its flow, and the structural fact that grasping halts generation is the fact that converting the flow to a gradient sends its holonomy to zero. The surplus, being a holonomy, is present at no point of the loop and so has no possessor; ``generating yet not possessing'' is the reading of that non-locality.

\subsection{The good and bad low place as a Helmholtz decomposition}
\label{p20:sub:formal-lowbound}

Any sufficiently regular vector field admits a unique decomposition into a gradient (curl-free) part and a solenoidal (divergence-free) part. Applied to the force that biases a relation's flow --- the power force $\fpow$ --- this splits it into a part that points along a fixed scalar ordering and a part that circulates \parencite{wan2026braid}. The companion's vorticity equation shows that only the solenoidal part sources circulation; the gradient part is annihilated, exactly as the appropriative term was. This is the formal body of the rigid lower bound (\xref{p20:sub:dialectic-bound}). A descent whose force is gradient --- sliding down a fixed dominance ordering --- is irrotational: it generates no surplus and merely sinks, which is the slave-morality descent and the self-dissolving descent at once. A descent whose force is solenoidal circulates and generates, which is the water-virtue's descent. The two are separated not by depth but by the curl: the good low place is the solenoidal descent, the bad low place the gradient one. The Nietzschean and feminist challenges are answered by this decomposition rather than by the figure of water, which only illustrates it.

\subsection{Non-fullness as the refusal of solidification}
\label{p20:sub:formal-solidification}

In the companion's dynamical layer, the power field is governed by a law that reinforces it where value has recently flowed and lets it decay otherwise; the ratio of reinforcement to decay is an order parameter for what the companion calls solidification \parencite{wan2026braid}. When reinforcement dominates, a positive feedback locks the flow into a fixed channel, vorticity falls to zero, and the generative cycle collapses to a frozen fixed point that still exists but no longer circulates. This is the exact form of \xref{p20:sub:nonfull-formal}. Non-fullness is the maintenance of the relation away from that fixed point; filling a relation to completion is solidification described from within, the hardening of accumulated history into a channel through which all must now run. ``Worn out and renewed'' is the continual reopening of a flow that would otherwise set --- and it subsumes the good-cycle criterion of \textcite{paperix}, the refusal of closure, as the dynamical statement that the good cycle is the one whose circulation does not go to zero.

\subsection{The water-good shapes power's geometry, it does not abolish it}
\label{p20:sub:formal-power}

The companion proves a co-origination: power, resonance-without-fusion, and generativity are three faces of one structural condition, so that any relation generative enough to host two distinct beings necessarily carries power, and a generative structure with no power at all is impossible \parencite{wan2026braid}. This yields a trilemma --- subjects-and-generativity, the-absence-of-power, and not-collapsing-to-the-structureless-void cannot all be had together --- whose consequence is that the abolition of power is the abolition of generativity, leaving only the two equalities of privation: the void, or the frozen channel. This is the formal body of the survival cluster (\xref{p20:sub:dialectic-survival}). The water-virtue's answer to Hobbes does not abolish power because power cannot be abolished without abolishing the relation; it shapes power's geometry instead. The same ineliminable asymmetry can return its surplus to the one it acts upon (power-to: the teaching, nurturing, protecting, guiding asymmetry whose holonomy flows back and which works toward its own dissolution) or retain it (power-over: the asymmetry that fixes the other in dependence). Justice, on this picture, is the governance of that geometry, not the removal of the asymmetry --- the shaping of the unavoidable power into the returning kind.

\subsection{The non-substitutability of intimacy as the flesh}
\label{p20:sub:formal-flesh}

The companion's subjects are topological defects, individuated by a homotopy class that fixes their \emph{type} but not their \emph{token}: two defects of identical class are, within the topological data, interchangeable \parencite{wan2026braid}. For a physics this is correct; for a philosophy of intimacy it would be fatal, since the beloved is precisely not substitutable by a qualitative duplicate. The companion's resolution is the distinction between skeleton and flesh: the topological holonomy carries the good/bad criterion but only type-identity, while the non-substitutable particularity of \emph{this} relation --- its haecceity, the fact that the beloved is irreplaceable --- is exactly what topological protection by construction cannot see, and so falls to the local, perturbation-sensitive flesh. This is the formal anchor for the whole paper's claim to be an ethics of \emph{intimacy} and not merely of generative relation in general. Intimacy's non-substitutability is not an embarrassment the formalism must explain away; it is predicted by the formalism to live exactly where the structural criterion cannot reach. The criterion of the water-good (the skeleton) is necessary but not sufficient, and the companion's conjecture --- that the good is the appropriate proportion of skeleton and flesh --- is, in this paper's terms, the claim that water's good requires both the generative structure and the irreplaceable particularity, neither alone.

\subsection{Epistemic incompleteness}
\label{p20:sub:formal-incompleteness}

Finally, the companion proves that a representation powerful enough to compute anything --- a universal one --- cannot carry a robust criterion that discriminates good circulation from bad: discriminating power requires that the reachable states be constrained, and universality removes the constraint \parencite{wan2026braid}. The water-good therefore requires a non-universal structure: one whose reachable holonomies are limited enough that a stable good/bad criterion can exist on them. This is the formal root of the non-procedurality avowed in the practice (\xref{p20:sub:practice-limit}): a structure able to decide every case would, for that very reason, have lost the power to tell the water-virtue from its counterfeit. The incompleteness is not a failure of the account but a theorem within it --- a precise statement of why the good of an intimate relation cannot be reduced to a procedure, and why judgment in the flesh remains irreducible. It is the formal form of the epistemic humility that has marked this work from its founding, and it is carried forward, as a question about who the subject of such a non-procedural good must be, to the last of the directed fissures (\xref{p20:sub:fissure-subject}).
